% GENERATED FILE. Edit canonical lessons, then run npm run generate:books.
% canonical-source-hash: fnv1a64:d099c32792f4dfe3
% canonical-lessons: KA-C441-sarati, KA-C441-manna-madu, KA-C441-desiya, KA-C441-nagadu, KA-C441-kard

\chapter{Paying Charges}
\label{ch:ka-paying-charges}
\hlchaptermodality{441}

\begin{chapteropening}
\textbf{By the end of this chapter:} \emph{I can say a turn in line, to waive, domestic, cash and a card.}

Complete the last lesson of chapter 441: I can say a turn in line, to waive, domestic, cash and a card.
\end{chapteropening}

\section[\texorpdfstring{\kn{ಸರತಿ} (sarati)}{ಸರತಿ (sarati)}]{\kn{ಸರತಿ} (sarati) — a turn in line; a queue}

\label{lesson:KA-C441-sarati}



\begin{quote}
Before the new one: say the Kannada for an embassy, then the Kannada for international.
\end{quote}

\subsection*{You'll want to know: \kn{ಸರತಿ}}
\textbf{\kn{ಸರತಿ}} — \emph{sarati} — \textquotedblleft{}a turn in line; a queue\textquotedblright{}.

\begin{quote}
\kn{ಬಿಲ್} \kn{ಕಟ್ಟಲು} \kn{ನಾವು} \kn{ಸರತಿಯಲ್ಲಿ} \kn{ನಿಂತೆವು}.
\end{quote}

\begin{quote}
\emph{bil kaṭṭalu nāvu saratiyalli nintevu.}
\end{quote}

\begin{quote}
We stood in the queue to pay the bill.
\end{quote}

\kn{ಸರತಿ} \kn{ಸಾಲು}, a queue; \kn{ನನ್ನ} \kn{ಸರತಿ}, my turn.

\begin{cousinweb}
\textbf{\kn{ಸರತಿ}} (\emph{sarati}) is a Kannada word for a turn in a line whose origin is not settled.
\end{cousinweb}

\begin{grammarlens}[title={What you've built}]
One more word for getting around.
\end{grammarlens}

\subsection*{Your turn}
\begin{itemize}
\raggedright
  \item \emph{Say it:} \emph{sarati}
  \item \emph{Say it:} \emph{sarati}, once more
  \item \emph{Say it:} \emph{bil kaṭṭalu nāvu saratiyalli nintevu.} — We stood in the queue to pay the bill.
  \item \emph{Recall:} say the Kannada for an embassy, then the Kannada for international, then say \emph{sarati} again
\end{itemize}

\begin{tcolorbox}[breakable,colback=teal!4,colframe=teal!35!black,title={Before you move on}]
What does \kn{ಸರತಿ} mean? (\textquotedblleft{}A turn in line; a queue\textquotedblright{}.) Say it once more.
\end{tcolorbox}
\section[\texorpdfstring{\kn{ಮನ್ನಾ} \kn{ಮಾಡು} (mannā māḍu)}{ಮನ್ನಾ ಮಾಡು (mannā māḍu)}]{\kn{ಮನ್ನಾ} \kn{ಮಾಡು} (mannā māḍu) — to waive; to write off}

\label{lesson:KA-C441-manna-madu}



\begin{quote}
Before the new one: say the Kannada for international, then the Kannada for a turn in line.
\end{quote}

\subsection*{You'll want to know: \kn{ಮನ್ನಾ} \kn{ಮಾಡು}}
\textbf{\kn{ಮನ್ನಾ} \kn{ಮಾಡು}} — \emph{mannā māḍu} — \textquotedblleft{}to waive; to write off\textquotedblright{}.

\begin{quote}
\kn{ಹೋಟೆಲ್} \kn{ಮಕ್ಕಳ} \kn{ಶುಲ್ಕವನ್ನು} \kn{ಮನ್ನಾ} \kn{ಮಾಡಿತು}.
\end{quote}

\begin{quote}
\emph{hōṭel makkaḷa śulkavannu mannā māḍitu.}
\end{quote}

\begin{quote}
The hotel waived the fee for the children.
\end{quote}

Past (I …): \kn{ಮನ್ನಾ} \kn{ಮಾಡಿದೆ}, \emph{mannā māḍide}.

Said of fees, fines and debts: \kn{ದಂಡ} \kn{ಮನ್ನಾ}, a fine waived.

\begin{cousinweb}
\textbf{\kn{ಮನ್ನಾ} \kn{ಮಾಡು}} (\emph{mannā māḍu}) joins \emph{mannā}, \textquotedblleft{}waived, remitted\textquotedblright{}, which came into Kannada through Hindustani and whose earlier source is not settled, and the native verb \kn{ಮಾಡು}, to do.
\end{cousinweb}

\begin{grammarlens}[title={What you've built}]
One more word for getting around.
\end{grammarlens}

\subsection*{Your turn}
\begin{itemize}
\raggedright
  \item \emph{Say it:} \emph{mannā māḍu}
  \item \emph{Say it:} \emph{mannā māḍu}, once more
  \item \emph{Say it:} \emph{hōṭel makkaḷa śulkavannu mannā māḍitu.} — The hotel waived the fee for the children.
  \item \emph{Recall:} say the Kannada for international, then the Kannada for a turn in line, then say \emph{mannā māḍu} again
\end{itemize}

\begin{tcolorbox}[breakable,colback=teal!4,colframe=teal!35!black,title={Before you move on}]
What does \kn{ಮನ್ನಾ} \kn{ಮಾಡು} mean? (\textquotedblleft{}To waive; to write off\textquotedblright{}.) Say it once more.
\end{tcolorbox}
\section[\texorpdfstring{\kn{ದೇಶೀಯ} (dēśīya)}{ದೇಶೀಯ (dēśīya)}]{\kn{ದೇಶೀಯ} (dēśīya) — domestic; of the country}

\label{lesson:KA-C441-desiya}



\begin{quote}
Before the new one: say the Kannada for a turn in line, then the Kannada for to waive.
\end{quote}

\subsection*{You'll want to know: \kn{ದೇಶೀಯ}}
\textbf{\kn{ದೇಶೀಯ}} — \emph{dēśīya} — \textquotedblleft{}domestic; of the country\textquotedblright{}.

\begin{quote}
\kn{ದೇಶೀಯ} \kn{ವಿಮಾನಗಳು} \kn{ಒಂದನೇ} \kn{ಟರ್ಮಿನಲ್ನಿಂದ} \kn{ಹೊರಡುತ್ತವೆ}.
\end{quote}

\begin{quote}
\emph{dēśīya vimānagaḷu ondanē ṭarminalninda horaḍuttave.}
\end{quote}

\begin{quote}
Domestic flights leave from terminal one.
\end{quote}

Its partner on airport signs is \kn{ಅಂತರರಾಷ್ಟ್ರೀಯ}, international.

\begin{cousinweb}
\textbf{\kn{ದೇಶೀಯ}} (\emph{dēśīya}) is borrowed from Sanskrit \emph{dēśīya}, \textquotedblleft{}of the country\textquotedblright{}, built on \emph{dēśa}, \textquotedblleft{}a region, a country\textquotedblright{} (CDIAL 6547), the source of \kn{ದೇಶ}.
\end{cousinweb}

\begin{grammarlens}[title={What you've built}]
One more word for getting around.
\end{grammarlens}

\subsection*{Your turn}
\begin{itemize}
\raggedright
  \item \emph{Say it:} \emph{dēśīya}
  \item \emph{Say it:} \emph{dēśīya}, once more
  \item \emph{Say it:} \emph{dēśīya vimānagaḷu ondanē ṭarminalninda horaḍuttave.} — Domestic flights leave from terminal one.
  \item \emph{Recall:} say the Kannada for a turn in line, then the Kannada for to waive, then say \emph{dēśīya} again
\end{itemize}

\begin{tcolorbox}[breakable,colback=teal!4,colframe=teal!35!black,title={Before you move on}]
What does \kn{ದೇಶೀಯ} mean? (\textquotedblleft{}Domestic; of the country\textquotedblright{}.) Say it once more.
\end{tcolorbox}
\section[\texorpdfstring{\kn{ನಗದು} (nagadu)}{ನಗದು (nagadu)}]{\kn{ನಗದು} (nagadu) — cash}

\label{lesson:KA-C441-nagadu}



\begin{quote}
Before the new one: say the Kannada for to waive, then the Kannada for domestic.
\end{quote}

\subsection*{You'll want to know: \kn{ನಗದು}}
\textbf{\kn{ನಗದು}} — \emph{nagadu} — \textquotedblleft{}cash\textquotedblright{}.

\begin{quote}
\kn{ಈ} \kn{ಅಂಗಡಿಯಲ್ಲಿ} \kn{ನಗದು} \kn{ಮಾತ್ರ} \kn{ತೆಗೆದುಕೊಳ್ಳುತ್ತಾರೆ}.
\end{quote}

\begin{quote}
\emph{ī aṅgaḍiyalli nagadu mātra tegedukoḷḷuttāre.}
\end{quote}

\begin{quote}
This shop takes cash only.
\end{quote}

\kn{ನಗದು} \kn{ಕೊಡು}, to pay in cash.

\begin{cousinweb}
\textbf{\kn{ನಗದು}} (\emph{nagadu}) is usually traced to Arabic \emph{naqd}, \textquotedblleft{}ready money\textquotedblright{}, which reached Kannada through Hindustani.
\end{cousinweb}

\begin{grammarlens}[title={What you've built}]
One more word for getting around.
\end{grammarlens}

\subsection*{Your turn}
\begin{itemize}
\raggedright
  \item \emph{Say it:} \emph{nagadu}
  \item \emph{Say it:} \emph{nagadu}, once more
  \item \emph{Say it:} \emph{ī aṅgaḍiyalli nagadu mātra tegedukoḷḷuttāre.} — This shop takes cash only.
  \item \emph{Recall:} say the Kannada for to waive, then the Kannada for domestic, then say \emph{nagadu} again
\end{itemize}

\begin{tcolorbox}[breakable,colback=teal!4,colframe=teal!35!black,title={Before you move on}]
What does \kn{ನಗದು} mean? (\textquotedblleft{}Cash\textquotedblright{}.) Say it once more.
\end{tcolorbox}
\section[\texorpdfstring{\kn{ಕಾರ್ಡ್} (kārḍ)}{ಕಾರ್ಡ್ (kārḍ)}]{\kn{ಕಾರ್ಡ್} (kārḍ) — a card; a bank card}

\label{lesson:KA-C441-kard}



\begin{quote}
Before the new one: say the Kannada for domestic, then the Kannada for cash.
\end{quote}

\subsection*{You'll want to know: \kn{ಕಾರ್ಡ್}}
\textbf{\kn{ಕಾರ್ಡ್}} — \emph{kārḍ} — \textquotedblleft{}a card; a bank card\textquotedblright{}.

\begin{quote}
\kn{ನಾನು} \kn{ಟಿಕೆಟ್ಗೆ} \kn{ಕಾರ್ಡ್} \kn{ಮೂಲಕ} \kn{ಪಾವತಿಸಿದೆ}.
\end{quote}

\begin{quote}
\emph{nānu ṭikeṭge kārḍ mūlaka pāvatiside.}
\end{quote}

\begin{quote}
I paid for the ticket by card.
\end{quote}

\kn{ಕಾರ್ಡ್} \kn{ಮೂಲಕ}, by card; \kn{ನಗದಾಗಿ}, in cash.

\begin{cousinweb}
\textbf{\kn{ಕಾರ್ಡ್}} (\emph{kārḍ}) is a loan from English \emph{card}.
\end{cousinweb}

\begin{grammarlens}[title={What you've built}]
One more word for getting around.
\end{grammarlens}

\subsection*{Your turn}
\begin{itemize}
\raggedright
  \item \emph{Say it:} \emph{kārḍ}
  \item \emph{Say it:} \emph{kārḍ}, once more
  \item \emph{Say it:} \emph{nānu ṭikeṭge kārḍ mūlaka pāvatiside.} — I paid for the ticket by card.
  \item \emph{Recall:} say the Kannada for domestic, then the Kannada for cash, then say \emph{kārḍ} again
\end{itemize}

\begin{tcolorbox}[breakable,colback=teal!4,colframe=teal!35!black,title={Before you move on}]
What does \kn{ಕಾರ್ಡ್} mean? (\textquotedblleft{}A card; a bank card\textquotedblright{}.) Say it once more.
\end{tcolorbox}
